% =====================================================
% 题型：正弦定理与三角恒等变形解答题 (q_tjw_20260606_24_sine_rule_angle_expression.tex)
% =====================================================
% =====================================================
% 难度：★★★ (难题)
% =====================================================
\newcommand{\qTriangleSineRuleExpressionStem}{%
  在 \(\triangle ABC\) 中，内角 \(A,B,C\) 所对的边分别为 \(a,b,c\)。\\
  已知 \(b\sin A = a\cos\left(B-\dfrac{\pi}{6}\right)\)。\\
  （I）求角 \(B\) 的大小；\\
  （II）设 \(a=2\)，\(c=3\)，求 \(b\) 和 \(\sin(2A-B)\) 的值。%
}

\newcommand{\qTriangleSineRuleExpressionAnswer}{%
  （I）\(B=\dfrac{\pi}{3}\)；（II）\(b=\sqrt7\)，\(\sin(2A-B)=\dfrac{3\sqrt3}{14}\)。%
}

\newcommand{\qTriangleSineRuleExpressionSolution}{%
  （I）由正弦定理
  \[
    \frac{a}{\sin A}=\frac{b}{\sin B},
  \]
  可得
  \[
    b\sin A = a\sin B.
  \]
  又已知
  \[
    b\sin A = a\cos\left(B-\frac{\pi}{6}\right),
  \]
  且 \(a>0\)，所以
  \[
    \sin B = \cos\left(B-\frac{\pi}{6}\right).
  \]

  展开右边：
  \[
    \cos\left(B-\frac{\pi}{6}\right)
    = \frac{\sqrt3}{2}\cos B+\frac{1}{2}\sin B.
  \]
  因此
  \[
    \sin B
    = \frac{\sqrt3}{2}\cos B+\frac{1}{2}\sin B,
    \qquad
    \sin B = \sqrt3\cos B.
  \]
  因为 \(B\in(0,\pi)\)，且右式要求 \(\cos B>0\)，所以 \(B\in\left(0,\frac{\pi}{2}\right)\)。
  故
  \[
    \tan B = \sqrt3,
    \qquad
    B=\frac{\pi}{3}.
  \]

  （II）由余弦定理得
  \[
    b^2
    = a^2+c^2-2ac\cos B
    = 2^2+3^2-2\cdot2\cdot3\cdot\frac12
    = 7.
  \]
  所以 \(b=\sqrt7\)。

  由正弦定理得
  \[
    \frac{\sin A}{a}=\frac{\sin B}{b},
    \qquad
    \sin A
    = \frac{a\sin B}{b}
    = \frac{2\cdot\frac{\sqrt3}{2}}{\sqrt7}
    = \frac{\sqrt3}{\sqrt7}.
  \]
  因为 \(a<c\)，所以 \(A<C\)，从而 \(A\) 为锐角，故
  \[
    \cos A=\sqrt{1-\sin^2 A}
    = \sqrt{1-\frac{3}{7}}
    = \frac{2}{\sqrt7}.
  \]

  于是
  \[
    \sin 2A
    = 2\sin A\cos A
    = \frac{4\sqrt3}{7},
    \qquad
    \cos 2A
    = \cos^2A-\sin^2A
    = \frac{1}{7}.
  \]
  因此
  \[
    \begin{aligned}
    \sin(2A-B)
      &= \sin2A\cos B-\cos2A\sin B \\
      &= \frac{4\sqrt3}{7}\cdot\frac12
       - \frac17\cdot\frac{\sqrt3}{2} \\
      &= \frac{3\sqrt3}{14}.
    \end{aligned}
  \]
}

\newcommand{\qTriangleSineRuleExpressionDiagnosis}{%
  （留白待收集学生作答后补充）%
}

\newcommand{\qTriangleSineRuleExpressionTeachingNotes}{%
  必讲候选。第（I）问是边角互化：先用正弦定理把 \(b\sin A\) 化为 \(a\sin B\)，再处理 \(\cos(B-\frac{\pi}{6})\)。第（II）问串联余弦定理、正弦定理和二倍角公式，适合作为本组综合收束题。%
}


% =====================================================
% 别名兼容：旧课次宏定义
% =====================================================
\newcommand{\qTJWZeroSixTwentyFourStem}{\qTriangleSineRuleExpressionStem}
\newcommand{\qTJWZeroSixTwentyFourAnswer}{\qTriangleSineRuleExpressionAnswer}
\newcommand{\qTJWZeroSixTwentyFourSolution}{\qTriangleSineRuleExpressionSolution}
\newcommand{\qTJWZeroSixTwentyFourDiagnosis}{\qTriangleSineRuleExpressionDiagnosis}
\newcommand{\qTJWZeroSixTwentyFourTeachingNotes}{\qTriangleSineRuleExpressionTeachingNotes}
